% ============================================================
%  SOS/A FORMELSAMMLUNG — Teil 7
%  Die harmonische Balance
% ============================================================

\bereich[cKap7]{Teil 7: Harmonische Balance — Grenzzyklen}

\begin{formelreihe}
  \parbox[t]{\linewidth}{%
    {\scriptsize\scshape\color{gray!65!black} Def.: Grenzzyklus}\\[0.03em]%
    \scriptsize
    Period.\,Lösung der Zustandsgl.\\
    (geschl.\,Trajektorie); in Nähe\\
    keine weitere period.\,Lösung.\\[0.06em]
    {\scriptsize\scshape\color{gray!65!black} Arten}\\[0.03em]%
    \textbf{Stabil:} Traj.\,konvergieren\\
    von innen \& außen auf GZ\\[0.03em]
    \textbf{Instabil:} Traj.\,divergieren\\
    vom GZ weg} &
  \parbox[t]{\linewidth}{%
    {\scriptsize\scshape\color{gray!65!black} Instabile Ruhelage $\Rightarrow$}\\[0.03em]%
    \scriptsize
    – Zustand $\to$ andere Ruhelage\\[0.02em]
    – Zustand $\to\infty$\\[0.02em]
    – Zustand $\to$ \textbf{Grenzzyklus}\\
    \quad(nichtlin.\,Dauerschwingung)\\[0.02em]
    – Chaotisches Verhalten\\[0.06em]
    {\scriptsize\scshape\color{gray!65!black} NL-Standardregelkreis}\\[0.03em]%
    \scriptsize
    – NL: durch Nullpunkt, symmetr.\\
    – $G(s)$: stabiles LZI-Teilsystem\\
    – Betrachte Grundschwingung} &
  \beispielblock{Einheitskreis-GZ}{%
    \dot{x}_1\!=\!-x_2\!+\!(1\!-\!x_1^2\!-\!x_2^2)x_1\\[0.04em]
    \dot{x}_2\!=\!\phantom{-}x_1\!+\!(1\!-\!x_1^2\!-\!x_2^2)x_2\\[0.07em]
    \Rightarrow\text{GZ: Einheitskreis}\\
    \text{(stabil, ampl.-unabhäng.)}} \\[0.3em]
\end{formelreihe}
\bereichende

\bereich[cKap7]{Beschreibungsfunktion \& Schwingungsbedingung}

\begin{formelreihe}
  \parbox[t]{\linewidth}{%
    {\scriptsize\scshape\color{gray!65!black} Grundschwingungsnäherung}\\[0.03em]%
    \scriptsize
    $v(t)\!\approx\!v_0\cos(\omega_0 t)$\\[0.03em]
    stat., symm.\,Kennlinie $\Rightarrow$\\
    Grundschw.\,$a(t)$ phasengleich\\
    mit $v(t)$\\[0.06em]
    {\scriptsize\scshape\color{gray!65!black} Beschreibungsfunktion $K_N(v_0)$}\\[0.03em]%
    $\displaystyle K_N(v_0)=\frac{a_1}{v_0}$\\[0.04em]
    {\scriptsize Fourier-Grundkoeff.\,von $a(t)$:}\\[0.02em]
    $\displaystyle a_1=\tfrac{2}{T_P}\!\int_0^{T_P}\!a(t)\cos\tfrac{2\pi t}{T_P}\,dt$\\[0.03em]
    \scriptsize NL $\to$ amplitudenabh.\,Verstärk.} &
  \parbox[t]{\linewidth}{%
    {\scriptsize\scshape\color{cKap7!80!black} Schwingungsbedingung}\\[0.03em]%
    $\displaystyle G(j\omega_0)=-\frac{1}{K_N(v_0)}$\\[0.06em]
    {\scriptsize\scshape\color{gray!65!black} Zwei-Ortskurven-Methode}\\[0.03em]%
    \scriptsize
    1.\;$G(j\omega)$-Ortskurve in kompl.\,Ebene\\[0.02em]
    2.\;Neg.\,inv.\,BF $-\tfrac{1}{K_N(v_0)}$ eintragen\\[0.02em]
    3.\;Schnittpunkt $=$ GZ-Kandidat\\[0.02em]
    \quad $v_0$: Beziff.\,der inv.-BF-Kurve\\[0.01em]
    \quad $\omega_0$: Beziff.\,der $G(j\omega)$-Kurve\\[0.04em]
    Periodendauer:\;$T_0=\tfrac{2\pi}{\omega_0}$\\[0.04em]
    Ampl.\,$v_0$ am Ausgang von $G(s)$} &
  \beispielblock{Schalter, $G\!=\!\tfrac{1}{(1+s)^3}$}{%
    K_N=\tfrac{4\cdot 1}{\pi v_0}\\[0.04em]
    -\tfrac{1}{K_N}=-\tfrac{\pi v_0}{4}\\[0.04em]
    \text{(neg.\,reelle Achse)}\\[0.05em]
    \Rightarrow\omega_0\!\approx\!1{,}7\\[0.02em]
    \Rightarrow T_0\!\approx\!3{,}7\\[0.02em]
    \Rightarrow v_0\!\approx\!0{,}16} \\[0.3em]
\end{formelreihe}
\bereichende

\bereich[cKap7]{Rezept: Dauerschwingung (GZ) bestimmen \& Amplituden}

\begin{rezept}
\rs{\textbf{Standardregelkreis bilden:} alle LZI-Blöcke zu
  $G(s)=V_1 V_2\,F(s)$ zusammenfassen (statische Verstärker $\cdot$ dynam.\ Teil), NL als Block $K_N(v_0)$ separat.}
\rs{$G(j\omega)$-Ortskurve (beziffert mit $\omega$) in kompl.\ Ebene eintragen.}
\rs{Neg.\ inv.\ BF $-\tfrac{1}{K_N(v_0)}$ eintragen (Lage s.\ BF-Tabelle: reelle Achse / Hysterese-Parallele).}
\rs{\textbf{Schwingungsbed.}\ $G(j\omega_0)=-\tfrac{1}{K_N(v_0)}$ = \textbf{Schnittpunkt} beider Kurven:\;
  $\omega_0$ an der $G$-Kurve ablesen,\; $v_0$ an der $-1/K_N$-Kurve ablesen.}
\rs{Periodendauer $T_0=\tfrac{2\pi}{\omega_0}$.}
\rs{\textbf{Stabilität:} $v_0$ leicht erhöhen ($+\varepsilon$). Liegt $-\tfrac{1}{K_N(v_0+\varepsilon)}$ \textbf{links} der $G(j\omega)$-Kurve (Durchlaufsinn wachsender $\omega$)?\; Ja $\Rightarrow$ stabiler GZ, Nein $\Rightarrow$ instabil.}
\end{rezept}

\miniexample{Zweipunkt-Schalter (Höhe $h$), $V_1\!=\!V_2\!=\!1$}{%
$F(j\omega)$ schneide die neg.\ reelle Achse bei $-1$; dort sei $\omega_0=0{,}65$ abgelesen.
Schalter: $K_N=\tfrac{4h}{\pi v_0}$, also $-\tfrac{1}{K_N}=-\tfrac{\pi v_0}{4h}$ (reell).
Schnittpunkt: $-\tfrac{\pi v_0}{4h}=-1\;\Rightarrow\;v_0=\tfrac{4h}{\pi}$.\quad
$T_0=\tfrac{2\pi}{0{,}65}\approx 9{,}7$. Schalter $\Rightarrow$ stets stabiler GZ.}

\bereichende

\bereich[cKap7]{Rezept: Signal-Amplituden im Regelkreis (Propagation)}

\noindent\footnotesize $v_0$ (HB) $=$ Amplitude am \textbf{NL-Eingang} (Arg.\ von $K_N$).\; Übrige Amplituden per \textbf{Zurück-/Vorwärtsrechnen} durch die Blöcke (je Grundschw., Frequenz $\omega_0$):\par\vspace{0.15em}
\begin{rezept}
\rs{Statischer Verstärker $V$:\; Ausgangsamplitude $=V\cdot$ Eingangsamplitude.}
\rs{Dynam.\ Block $F$:\; Ausgangsamplitude $=|F(j\omega_0)|\cdot$ Eingangsamplitude.\;
  $|F(j\omega_0)|$ aus Ortskurve (Betrag des Zeigers) bzw.\ $=\tfrac{|G(j\omega_0)|}{V_1V_2}$.}
\rs{Kette $b\!\to\!V_1\!\to\!u\!\to\!F\!\to\!x\!\to\!V_2\!\to\!v$:\;
  $x_0=\tfrac{v_0}{V_2}$,\;\; $u_0=\tfrac{x_0}{|F(j\omega_0)|}$,\;\; $b_0=\tfrac{u_0}{V_1}$.}
\end{rezept}

\miniexample{$V_1\!=\!V_2\!=\!2$, abgelesen $v_0=5$, $|F(j\omega_0)|=1$}{%
$x_0=\tfrac{v_0}{V_2}=2{,}5$;\quad $u_0=\tfrac{x_0}{|F(j\omega_0)|}=2{,}5$;\quad $b_0=\tfrac{u_0}{V_1}=1{,}25$.}

\bereichende

\bereich[cKap7]{Rezept: Entwurf — Verstärkung / Schalthöhe für Soll-Amplitude}

\begin{rezept}
\rs{Soll-Amplitude (z.\,B.\ $x_0$) über die Blockkette in $v_0$ am NL-Eingang umrechnen (s.\ Propagation).}
\rs{$K_N(v_0)$ aus BF bestimmen (Graph ablesen oder Formel einsetzen).}
\rs{Schwingungsbed.\ am Schnittpunkt erzwingen: freie Verstärkung so wählen, dass $G(j\omega_0)=-\tfrac{1}{K_N(v_0)}$.
  Schneidet $F$ die neg.\ reelle Achse beim Wert $r$, gilt $V_1V_2\,r=-\tfrac{1}{K_N(v_0)}$ $\Rightarrow$ freie Verstärkung lösen.}
\end{rezept}

\miniexample{Dreipunkt-Schalter ($h\!=\!d\!=\!1$), $V_2\!=\!1$: welches $V_1$ für $x_0=4$?}{%
$V_2=1\Rightarrow v_0=x_0=4$. BF-Graph: $K_N(4)=0{,}3$ $\Rightarrow -\tfrac{1}{K_N}=-3{,}33$.
$F$ schneidet neg.\ reelle Achse bei $-1$ $\Rightarrow V_1\cdot(-1)=-3{,}33$ $\Rightarrow \boxed{V_1=3{,}33}$.}

\fallstrick{%
\textbullet\ $v_0$ ist Amplitude am \textbf{NL-Eingang}, nicht am Ausgang $x$ — Blöcke umrechnen!\quad
\textbullet\ \textbf{Zweipunkt-Schalter:} neg.-inv.\ BF deckt die \emph{ganze} neg.\ reelle Achse ab $\Rightarrow$ $\omega_0$ (und $T_0$) liegt fest durch $G$; $h$ ändert nur $v_0$, \textbf{nicht} $T_0$.\quad
\textbullet\ \textbf{Dreipunkt/Begrenzer/Totzone:} $K_N=0$ für $v_0<d$; neg.-inv.\ BF beginnt erst bei endlichem Wert $\Rightarrow$ evtl.\ \textbf{kein Schnittpunkt} $\Rightarrow$ HB sagt keine Dauerschwingung.\quad
\textbullet\ \textbf{Kein Schnittpunkt $\not\Rightarrow$ kein GZ!} HB ist Näherung, exakt nur bei ausgeprägter TP-Wirkung von $G$ (Oberwellen gedämpft).\quad
\textbullet\ \textbf{Hysterese:} $K_N\in\mathbb{C}$ $\Rightarrow$ neg.-inv.\ BF \emph{nicht} auf der reellen Achse ($\operatorname{Im}=-\tfrac{d\pi}{4h}=$ const).}

\bereichende

\bereich[cKap7]{Beschreibungsfunktionen — Übersicht}

\noindent\footnotesize\textit{\textbf{ohne Gedächtnis} (stat., symm.): $K_N\in\mathbb{R}$ $\Rightarrow$ neg.\,inv.\,BF auf neg.\,reeller Achse.\quad\textbf{mit Gedächtnis} (Hyst.): $K_N\in\mathbb{C}$, $\operatorname{Im}\!\left(-\tfrac{1}{K_N}\right)=-\tfrac{d\pi}{4h}=\mathrm{konst.}$}\par\vspace{0.12em}
\noindent\footnotesize\setlength{\tabcolsep}{2pt}%
{\renewcommand{\arraystretch}{1.15}%
\begin{tabular}{@{}p{0.16\linewidth}@{\hspace{2pt}}p{0.50\linewidth}@{\hspace{3pt}}p{0.30\linewidth}@{}}
  \scriptsize\scshape\color{gray!65!black}Nichtlinearität &
  \scriptsize\scshape\color{gray!65!black}Beschreibungsfunktion $K_N(v_0)$ &
  \scriptsize\scshape\color{gray!65!black}Neg.\,inv.\,BF\;$-1/K_N$ \\[0.1em]
  Schalter &
  $\tfrac{4h}{\pi v_0}$ &
  $-\tfrac{\pi v_0}{4h}$ (reell)\\[0.2em]
  Schalter m.\,Hyst. &
  $\tfrac{4h}{\pi v_0}\sqrt{1\!-\!\tfrac{d^2}{v_0^2}}-j\tfrac{4dh}{\pi v_0^2}$ &
  $-\tfrac{\pi v_0}{4h}\sqrt{1\!-\!\tfrac{d^2}{v_0^2}}-j\tfrac{d\pi}{4h}$\\[0.04em]
  & \scriptsize\itshape $v_0\!\geq\!d$ (sonst kommt $v_0\!<\!d$ nicht vor) &
    \scriptsize\itshape $\operatorname{Im}=-\tfrac{d\pi}{4h}=\mathrm{konst.}$\\[0.15em]
  Dreipunktschalter &
  $\begin{cases}0 & v_0\!<\!d\\\tfrac{4h\sqrt{v_0^2-d^2}}{\pi v_0^2} & v_0\!\geq\!d\end{cases}$ &
  reell ($v_0\!\geq\!d$)\\[0.5em]
  Begrenzer &
  $\begin{cases}\tfrac{h}{d} & v_0\!<\!d\\[0.1em]\tfrac{2h}{d\pi}\!\left(\arcsin\tfrac{d}{v_0}\!+\!\tfrac{d}{v_0}\sqrt{1\!-\!\tfrac{d^2}{v_0^2}}\right) & v_0\!\geq\!d\end{cases}$ &
  reell\\[0.65em]
  Totzone &
  $\begin{cases}0 & v_0\!<\!d\\[0.1em]m\!-\!\tfrac{2m}{\pi}\!\left(\arcsin\tfrac{d}{v_0}\!+\!\tfrac{d}{v_0}\sqrt{1\!-\!\tfrac{d^2}{v_0^2}}\right) & v_0\!\geq\!d\end{cases}$ &
  reell\\
\end{tabular}}\par
\vspace{0.15em}
\bereichende

\bereich[cKap7]{Stabilität von Grenzzyklen}

\begin{formelreihe}
  \parbox[t]{\linewidth}{%
    {\scriptsize\scshape\color{gray!65!black} Stabilitätsprüfung am Schnittpunkt}\\[0.03em]%
    \scriptsize
    Am Schnittpunkt gilt $G(j\omega_0)=-\tfrac{1}{K_N(v_0)}$\\[0.04em]
    Erhöhe $v_0$ leicht ($v_0\!+\!\varepsilon$):\\[0.03em]
    Liegt $-\tfrac{1}{K_N(v_0\!+\!\varepsilon)}$\\
    \quad\textbf{links} der $G(j\omega)$-Ortskurve?\\[0.04em]
    $\Rightarrow$ \textbf{Ja}: stabiler GZ\;$\bigcirc$\\[0.03em]
    $\Rightarrow$ \textbf{Nein}: instabiler GZ\;$\triangle$} &
  \parbox[t]{\linewidth}{%
    {\scriptsize\scshape\color{gray!65!black} Gültigkeit der Methode}\\[0.03em]%
    \scriptsize
    exakt nur bei ausgepr.\,TP-Wirkung\\
    von $G(s)$ (Oberwellen gedämpft)\\[0.04em]
    \textbf{Kein Schnittpunkt} $\not\Rightarrow$ kein GZ!\\[0.03em]
    Bsp.: PT1\,+\,Schalter\,m.\,Hyst.\\
    $\Rightarrow$ GZ trotzdem\\[0.04em]
    ohne Schnittpunkt: keine\\
    Aussage über Schwingung.} &
  \beispielblock{Vorgehen kompakt}{%
    \text{1. }G(j\omega)\text{-Ortskurve}\\[0.03em]
    \text{2. }-\tfrac{1}{K_N(v_0)}\text{-Kurve}\\[0.03em]
    \text{3. Schnittpkt.\!}\to\!v_0,\omega_0\\[0.03em]
    \text{4. }T_0=\tfrac{2\pi}{\omega_0}\\[0.04em]
    \text{5. }v_0{\uparrow}\!:\text{ links?}\Rightarrow\!\text{stab.}} \\[0.3em]
\end{formelreihe}
\bereichende
